\chapter{Introduction}
\label{chap:intro}

Modern epidemiological and health research increasingly uses electronic health records (EHRs) and other large observational databases. These data were assembled for administration or care rather than research \citep{goldstein2016, hubbard2015}. They provide large cohorts at low cost, but many recorded variables are error-prone. For a fixed measurement-error mechanism, the bias this induces in a regression coefficient generally persists as the sample size grows, while the sampling error decreases \citep{carroll2006measurement}. An analysis of a large cohort can therefore produce a narrow confidence interval around a biased coefficient estimate.

Where a reference value can be recovered, a practical response is to audit a probability subsample. For each selected record, a reference value is established by reviewing the source record or by using a reference assay where one exists, while the error-prone value remains available for the full cohort. This is a two-phase design in the sense of \citet{neyman1938contribution}: phase one holds the error-prone data for everyone, and phase two holds the reference values for the audited records. Validation of this kind has been used to correct measurement error in EHR analyses \citep{giganti2020accounting, shepherd2023multiwave, amorim2021two}. Auditing only a subsample reduces the cost. Phase-one information can then be used to choose sampling probabilities that improve the precision of a target coefficient \citep{amorim2021two, shepherd2023multiwave}. The goal is to combine both phases into a regression coefficient on the reference scale, with a confidence interval that has correct coverage.

One way to use both phases is imputation. A model for the reference measurement is fitted on the validated rows and used to complete the unvalidated ones, and the completed data are analysed by the rules of multiple imputation \citep{Rubin1987}. The completed data can support several analysis models if the imputation model includes the variables and relationships required by those analyses. Imputation can use automated methods, including deep generative models developed to approximate distributions of tabular data \citep{yoon2018gain, kotelnikov2023tabddpm}. This thesis asks whether these methods give valid coefficient estimates and confidence intervals in two-phase validation studies.

A recognised caution is that automated multiple imputation can distort the relationships an analysis targets: adding auxiliary variables that are colliders, or otherwise tied to the reason data are missing, can open non-causal paths and increase rather than reduce bias \citep{thoemmes2014cautious}. This risk arises when the missingness mechanism is unknown and the analyst must judge which auxiliaries are safe to include. It does not apply to the designs studied here, because phase-two validation is drawn by a known probability sampling design on the phase-one data, so the reference measurement is missing at random given the full phase-one information. The variables that determine selection are known and observed, conditioning on them is the correct and complete adjustment, and the collider bias the caution describes cannot arise. The concern applies to general missing-data imputation, where missingness at random is assumed rather than guaranteed by design.

Imputation can be judged in two ways. One asks how closely the imputed values recover the withheld ones. The other asks whether the downstream analysis is valid, and is measured by the bias of the target coefficient and the coverage of its confidence interval \citep{naf2023imputation, Rubin1987, meng1994multiple}. A method can do well on the first and still fail the second, and the second is the criterion this thesis adopts.

Informative phase-two sampling can break that criterion. When the audit sample is drawn with unequal probabilities and the variables that drove the selection are withheld, the validated rows follow the sampled distribution of the measurement rather than its population distribution. An imputation model fitted on those rows without the design information reproduces the sampled relationship and biases the coefficient. The released sampling weight then has to enter the correction, through either the imputation model or the analysis \citep{lumley2011complex}.

Uncertainty in the fitted imputation model can break the same criterion in a separate way. The imputation model is estimated from data, so its fitted parameters carry uncertainty, and a valid interval must propagate that uncertainty. A single-fit generative procedure does not. One model is fitted, and the completed datasets are then produced by re-sampling that fixed model's noise, so they vary through the sampling draws but not through the fitted parameters. Rubin's between-imputation variance is computed across the completed datasets, so it reflects the sampling draws but leaves out the parameter uncertainty. The reported interval is then too narrow, and this can happen even when the coefficient itself is well estimated.

The thesis develops three imputation methods for these problems. SISL (sample imputation via sample likelihood) addresses informative phase-two sampling when only the design weight is released. It applies the sample-distribution identity of \citet{pfeffermann1999parametric} to the imputation likelihood and estimates the required inclusion function from the released weights and observed variables. SISL removes most of the bias from ignoring the informative selection. Its sandwich and bootstrap variants obtain near-nominal coverage by representing uncertainty in the fitted sample likelihood. The main result is negative. The suppressed stratifying variables prevent direct reconstruction of the original sampling rule, so SISL fits an inclusion surface. Weighted-likelihood imputation uses the realised weights directly, while design-based raking estimates the analysis coefficient through calibration. Both have lower bias than SISL in Simulation Design~I. The efficiency advantage established for sample-likelihood analysis does not carry over to this imputation comparison.

SICG (sample imputation via conditional generative adversarial networks) adapts a tabular generative adversarial network to impute the phase-two variables conditional on the observed record. Its coefficient estimates have low bias relative to most of the other imputation methods. Its intervals often under-cover because the retained uncertainty procedure does not supply enough variation across fitted models. Repeated draws with Monte Carlo dropout vary the masks within one trained generator, but they produce much less between-imputation variance than bootstrap retraining. In many settings, the resulting Rubin standard errors are smaller than the empirical sampling variation. The SICG results show that low-bias coefficient estimates can coexist with inadequate propagation of fitted-model uncertainty in a single-fit procedure.

SIRD (sample imputation via conditional residual denoising diffusion networks) combines a non-adversarial denoising objective with a diffusion path referenced to the phase-one proxy. For continuous variables, the path moves the validated value towards the paired proxy and learns the measurement-error residual. For categorical variables, proxy-to-truth transitions define the terminal distribution of a multinomial diffusion. The retained construction was cheap enough to refit on bootstrap resamples, and those fits carry fitted-model variation into the completed datasets. The full SIRD procedure gives low coefficient bias and improves interval coverage relative to SICG across the main coefficient blocks.

The methods are evaluated across three simulation designs and the Vanderbilt Comprehensive Care Clinic (VCCC) case study. The evaluation uses the bias and root mean squared error of the target coefficient estimate, its empirical and mean reported standard errors, its nominal coverage, and two oracle coverages that show how coverage changes after replacing either the reported standard error or the average simulation bias. When the stratifying variables that complete the sampling strata are suppressed, the methods that use the released weights outperform the general-purpose imputation methods. When the selection-related phase-one variables are observed, SICG and SIRD give low coefficient bias relative to the other imputation methods, and SIRD's bootstrap fits improve the representation of fitted-model uncertainty. Design-based raking generally has the lowest bias when its calibration constraints hold, although it can be badly biased under the tail-enrichment designs. Nearly all the imputation methods under-cover the interaction terms. This common failure is consistent with a limitation in preserving the product terms and is not specific to SICG or SIRD.

The thesis is organised as follows. Chapter~\ref{chap:background} sets out measurement error, the two-phase design, design-based and model-based estimation, and the variance decomposition of multiple imputation. Chapter~\ref{chap:ImpReview} reviews the comparator methods and the statistical and machine-learning work on tabular imputation, and states the fitted-parameter uncertainty problem. Chapter~\ref{chap:SimDesigns} defines the evaluation protocol: Simulation Design~I isolates informative sampling with suppressed selection variables, Simulation Design~II uses Gaussian and Binomial analysis models, Simulation Design~III uses a Cox model, and the chapter also specifies the VCCC case study and the reporting metrics. Chapters~\ref{chap:SISL}, \ref{chap:SICG} and~\ref{chap:SIRD} develop SISL, SICG and SIRD in turn, and Chapter~\ref{chap:conclusion} draws the results together, states their limits, and gives practical guidance. The proofs, tuning details, and full per-design results are collected in the appendices.
